\documentclass{article}
\PassOptionsToPackage{numbers,compress}{natbib}
\usepackage[final]{neurips}
\usepackage[utf8]{inputenc}
\usepackage[T1]{fontenc}
\usepackage{hyperref}
\usepackage{url}
\usepackage{booktabs}
\usepackage{amsfonts}
\usepackage{amsmath}
\usepackage{nicefrac}
\usepackage{microtype}
\usepackage{xcolor}
\usepackage{graphicx}

\title{VAE vs. DDPM for Image Generation on Fashion-MNIST}

\author{%
  Yuye Huang\\
  Department of Electrical and Computer Engineering\\
  University of Toronto\\
  \texttt{yuye.huang@mail.utoronto.ca} \\
  % more authors
  \And
  Jenny You\\
  Department of Mechanical and Industrial Engineering\\
  University of Toronto\\
  \texttt{Jenny.You@mail.utoronto.ca} \\
  \And
  Jasmine Shao\\
  Department of Mechanical and Industrial Engineering\\
  University of Toronto\\
  \texttt{jasmine.shao@mail.utoronto.ca} \\
}

\begin{document}
\maketitle

\begin{abstract}
This project compares a convolutional Variational Autoencoder (VAE) and a Denoising Diffusion Probabilistic Model (DDPM) for image generation on Fashion-MNIST. Both models were trained on the same dataset and evaluated with a shared pipeline using generated image samples, classifier-based feature metrics, class-distribution statistics, and measured computational cost. The DDPM produced sharper and more recognizable samples and achieved substantially lower classifier-feature FID and KID than the VAE. It also generated a more balanced class distribution. In contrast, the VAE required much less training time and generated samples almost immediately through a single decoder pass. The results show a clear trade-off between generation quality and computational efficiency for the two generative approaches.
\end{abstract}

\section*{Assentation of Teamwork}
All group members contributed to problem framing, implementation planning, experimental design, and result discussion. Jenny primarily worked on the VAE implementation, training pipeline, and latent-space experiments. Yuye primarily worked on the DDPM implementation, diffusion training and sampling pipeline, and denoising visualizations. Jasmine primarily handled the shared evaluatioGenAIn pipeline, quantitative metric analysis, result comparison, and report drafting.

\section{Introduction}
Generative models learn patterns from observed data and use them to create new samples. In this project, we compare two generative approaches covered in the course: Variational Autoencoders (VAEs) and Denoising Diffusion Probabilistic Models (DDPMs). A VAE represents an image through a lower-dimensional latent variable and generates new images by sampling and decoding from this latent space [1]. A DDPM instead starts from Gaussian noise and generates an image through a sequence of learned denoising steps [2].

The goal of this project is to compare these two approaches under the same Fashion-MNIST experimental setting. Fashion-MNIST contains 60,000 training images and 10,000 test images of 28$\times$28 grayscale clothing items from 10 classes [2]. Rather than selecting a single ``best'' model, we focus on their practical trade-offs. We compare generated sample quality, distribution similarity, class diversity, generation behavior, training stability, and computational cost. This allows us to study both how the two models generate images and what is gained or lost when using a more computationally expensive diffusion process.

\section{Preliminaries and Problem Formulation}
To compare these two approaches, we first briefly review their underlying formulations.

Let $x \sim p_{\mathrm{data}}(x)$ denote an image from the Fashion-MNIST data distribution. The general objective is to train a generative model that can produce new samples with a distribution close to $p_{\mathrm{data}}$. 

For the VAE, the encoder estimates a Gaussian posterior $q_\phi(z\mid x)$ with mean $\mu$ and log variance $\log\sigma^2$. A latent variable is sampled using the reparameterization trick, $z=\mu+\sigma\epsilon$, where $\epsilon\sim\mathcal{N}(0,I)$, and the decoder models $p_\theta(x\mid z)$. Training minimizes a reconstruction term together with the KL divergence between the approximate posterior and the standard Gaussian prior. The reconstruction term encourages image fidelity, while the KL term regularizes the latent space toward a smooth and sampleable Gaussian prior.

For the DDPM, a forward diffusion process gradually adds Gaussian noise to a clean image. At timestep $t$, the noisy sample can be written as
\begin{equation}
 x_t=\sqrt{\bar{\alpha}_t}x_0+\sqrt{1-\bar{\alpha}_t}\epsilon,
 \qquad \epsilon\sim\mathcal{N}(0,I).
\end{equation}
The neural network is trained to predict the added noise from $x_t$ and $t$. During generation, the model begins from random Gaussian noise and repeatedly applies the learned reverse denoising process until an image is produced. Unlike the VAE’s direct latent-space decoding, DDPM generation requires multiple sequential denoising steps during sampling.

\section{Design}

Both models are trained and evaluated under a shared experimental setting. Fashion-MNIST is used for all experiments, and each trained model generates 5,000 images for evaluation. Since the DDPM operates in a normalized diffusion space, its inputs are mapped to $[-1,1]$ during training and mapped back to $[0,1]$ before evaluation, whereas the VAE is trained directly on images in $[0,1]$.

To evaluate both generative models under the same criteria, we construct a shared evaluation pipeline based on a separately trained Fashion-MNIST classifier. The classifier achieves 91.34\% test accuracy and provides both class probabilities and 128-dimensional intermediate feature vectors.

The classifier outputs are used to evaluate generated-image recognizability, class distribution, classifier-feature FID/KID, and nearest-neighbor distances. Since the classifier is independent of both generative models, the same evaluation procedure is applied to all generated samples, providing a consistent evaluation framework for comparison.

\section{Methodology}
\textbf{VAE.} The implemented VAE uses a convolutional encoder with two stride-2 convolution layers, followed by a fully connected layer that outputs a 32-dimensional latent mean and log variance. After reparameterization, the decoder uses fully connected layers and two transposed convolutions to recover a 28$\times$28 image. The loss is
\begin{equation}
 \mathcal{L}_{\mathrm{VAE}}=\mathcal{L}_{\mathrm{BCE}}+\beta D_{\mathrm{KL}}(q_\phi(z\mid x)\|p(z)),
\end{equation}
with $\beta=1$. The model is trained with AdamW, cosine learning-rate decay, and gradient clipping.

\textbf{DDPM.} The implemented DDPM uses 100 diffusion steps with a cosine noise schedule for both training and sampling. Its noise-prediction network is a small U-Net with residual time-conditioned blocks, skip connections, sinusoidal timestep embeddings, and a self-attention block at the bottleneck. At every training iteration, a random timestep is selected, Gaussian noise is added to the input, and the network predicts that noise. The objective is mean squared error between the predicted and true noise. The DDPM is trained with AdamW, cosine learning-rate decay, and gradient clipping. An exponential moving average (EMA) of the model parameters with decay 0.995 is used for final sampling.

\textbf{Evaluation.}
The shared evaluation pipeline compares the two models from three complementary perspectives: generation quality, sample diversity, and computational efficiency. Generation quality is evaluated using classifier-feature FID [4], KID [5], classifier confidence, and the proportion of generated samples with confidence of at least 0.8. Since FID and KID are computed from the feature space of our Fashion-MNIST classifier rather than ImageNet Inception features, they are reported only for comparison between the two models within this project. Sample diversity is measured using class coverage and normalized class-distribution entropy. Computational efficiency is evaluated using training time, sampling time, trainable parameter count, and peak GPU memory. Sobel edge sharpness and nearest-neighbor distance are additionally reported as supporting diagnostics.

\section{Numerical Experiments}
\subsection{Experimental Setup}
The report-profile experiment uses the full 60,000-image Fashion-MNIST training set and 10,000-image test set with random seed 42. Training and evaluation were run in PyTorch 2.11.0 with CUDA 12.8 on an NVIDIA GeForce RTX 5070 Ti. Table~\ref{tab:config} summarizes the main model settings. The VAE is trained for 25 epochs and the DDPM for 35 epochs. Both use a batch size of 1024 in the recorded report run.

\begin{table}[t]
\caption{Main experimental configuration.}
\label{tab:config}
\centering
\small
\begin{tabular}{lcc}
\toprule
Setting & VAE & DDPM \\
\midrule
Epochs & 25 & 35 \\
Batch size & 1024 & 1024 \\
Learning rate & $2\times10^{-3}$ & $2\times10^{-4}$ \\
Latent dimension / diffusion steps & 32 & 100 \\
Trainable parameters & 1.70M & 5.22M \\
\bottomrule
\end{tabular}
\end{table}

\subsection{Qualitative Comparison}
Figure~\ref{fig:samples} compares real Fashion-MNIST samples with samples generated by both models. Both models learn recognizable clothing shapes, but the VAE samples are generally blurrier and some outlines are less distinct. The DDPM samples have clearer edges and more recognizable object structure. This visual difference is consistent with the quantitative results reported below.

\begin{figure}[t]
  \centering
  \includegraphics[width=0.96\linewidth]{figures/real_vae_ddpm_comparison.png}
  \caption{Real Fashion-MNIST images and generated samples from the VAE and DDPM.}
  \label{fig:samples}
\end{figure}

The models also show different generation behavior. Figure~\ref{fig:generation} shows a VAE interpolation between the latent means of two test images and a DDPM denoising trajectory. The VAE changes the decoded output smoothly as the latent representation is interpolated, while the DDPM gradually forms a recognizable image from random noise over repeated reverse steps.

\begin{figure}[t]
\centering
\begin{minipage}{0.48\linewidth}
\centering
\includegraphics[width=\linewidth]{figures/vae_latent_interpolation.png}
\small (a) VAE latent interpolation
\end{minipage}
\hfill
\begin{minipage}{0.48\linewidth}
\centering
\includegraphics[width=\linewidth]{figures/ddpm_denoising_trajectory.png}
\small (b) DDPM denoising trajectory
\end{minipage}
\caption{Different generation structures of the two models.}
\label{fig:generation}
\end{figure}

\subsection{Quantitative Results}
Table~\ref{tab:results} gives the main comparison using 5,000 generated samples from each model. The DDPM achieves a much lower classifier-feature FID (7.66 versus 51.12) and KID (0.326 versus 1.874), indicating that its generated feature distribution is closer to the real reference set. The classifier also assigns higher average confidence to DDPM samples, and 65.6\% of DDPM samples reach at least 0.8 confidence compared with 41.0\% for the VAE.

\begin{table}[t]
\caption{Main quantitative comparison. Higher is better for confidence, recognizability, and entropy; lower is better for FID, KID, and time.}
\label{tab:results}
\centering
\small
\begin{tabular}{lrr}
\toprule
Metric & VAE & DDPM \\
\midrule
Classifier-feature FID $\downarrow$ & 51.12 & \textbf{7.66} \\
Classifier-feature KID $\downarrow$ & 1.874 & \textbf{0.326} \\
Mean classifier confidence $\uparrow$ & 0.704 & \textbf{0.832} \\
Confidence $\geq 0.8$ $\uparrow$ & 41.0\% & \textbf{65.6\%} \\
Class coverage & 10/10 & 10/10 \\
Normalized class entropy $\uparrow$ & 0.931 & \textbf{0.991} \\
Training time $\downarrow$ & \textbf{0.37 min} & 4.76 min \\
Sampling time $\downarrow$ & \textbf{0.0075 ms/img} & 8.02 ms/img \\
\bottomrule
\end{tabular}
\end{table}

Both models cover all 10 predicted Fashion-MNIST classes, but their distributions are not equally balanced. The VAE generates a relatively high proportion of T-shirt/top and Pullover samples while Coat and Sandal are underrepresented. Its normalized class entropy is 0.931. The DDPM distribution is closer to balanced across classes and reaches an entropy of 0.991. Both training procedures completed without non-finite losses. The final VAE test negative ELBO is 238.20, consisting of 221.98 reconstruction loss and 16.22 KL loss, while the DDPM noise-prediction MSE decreases to 0.0696. These loss values are not directly comparable because the two models optimize different objectives.

The largest cost difference appears during sampling. VAE sampling takes about 0.0075 ms per image in the recorded batched experiment, whereas DDPM sampling takes about 8.02 ms per image. This is expected from the implementations: the VAE generates an image with a single decoder pass, while the DDPM performs 100 sequential denoising steps. Training time also increases from 0.37 minutes for the VAE to 4.76 minutes for the DDPM. Peak GPU memory was 558 MB for the VAE and 6275 MB for the DDPM. A feature-space nearest-neighbor screen on 128 generated samples found 0\% of VAE samples and 0.78\% of DDPM samples below a strict threshold defined by the 1st percentile of real-image nearest-neighbor distances. We therefore did not observe strong evidence of direct training-set memorization, although this screen is not a proof that memorization is absent.

\section{Discussion}
The results show that the DDPM performs better on most of the measured generation-quality indicators in our implementation. Its sample grid is visually sharper, its classifier-feature FID and KID are substantially lower, and generated images receive higher classifier confidence. The class-distribution results also show why class coverage alone is not enough to describe diversity. Although both models generated all 10 predicted classes, the VAE distribution is noticeably more concentrated in several classes, while the DDPM is more balanced.

The main advantage of the VAE is efficiency. It uses fewer trainable parameters and requires much less measured training time and GPU memory. Its direct latent-space sampling is also over three orders of magnitude faster per image than the 100-step DDPM sampling procedure in this experiment. The two models therefore show a clear quality--efficiency trade-off rather than one model being better in every aspect.

There are several limitations to this comparison. The reported experiment uses one random seed, so it does not measure run-to-run variation. The DDPM is also larger than the VAE (5.22M versus 1.70M trainable parameters), so the experiment is a comparison of our two implemented systems rather than a parameter-matched architecture study. Finally, FID, KID, recognizability, and class-distribution measurements depend on our Fashion-MNIST classifier. Its 91.34\% test accuracy makes it useful for evaluation, but these classifier-based metrics should still be interpreted together with the visual samples.

\section{Conclusions}
We implemented and evaluated a convolutional VAE and a DDPM on Fashion-MNIST under a shared experimental pipeline. Both models generated recognizable samples, but the DDPM produced sharper images, better classifier-feature distribution scores, and a more balanced predicted class distribution. The VAE was much more computationally efficient, especially during sampling. Future work could repeat the experiment across multiple random seeds and compare models with more closely matched parameter counts or fewer DDPM sampling steps.

\section*{References}
\small
[1] D. P. Kingma and M. Welling. ``Auto-Encoding Variational Bayes.'' arXiv:1312.6114, 2013.

[2] J. Ho, A. Jain, and P. Abbeel. ``Denoising Diffusion Probabilistic Models.'' arXiv:2006.11239, 2020.

[3] H. Xiao, K. Rasul, and R. Vollgraf. ``Fashion-MNIST: A Novel Image Dataset for Benchmarking Machine Learning Algorithms.'' arXiv:1708.07747, 2017.

[4] M. Heusel, H. Ramsauer, T. Unterthiner, B. Nessler, and S. Hochreiter. ``GANs Trained by a Two Time-Scale Update Rule Converge to a Local Nash Equilibrium.'' Advances in Neural Information Processing Systems, 2017.

[5] M. Bi\'nkowski, D. J. Sutherland, M. Arbel, and A. Gretton. ``Demystifying MMD GANs.'' arXiv:1801.01401, 2018.

\section*{Appendix}
The following repository outputs are included as supplementary evidence for the main comparison.

\begin{figure}[h]
  \centering
  \includegraphics[width=0.82\linewidth]{figures/class_distribution_comparison.png}
  \caption{Predicted class distribution of 5,000 generated samples from each model. The dashed line marks a balanced 0.1 proportion.}
\end{figure}

\begin{figure}[h]
\centering
\begin{minipage}{0.48\linewidth}
\centering
\includegraphics[width=\linewidth]{figures/vae_total_loss.png}
\small (a) VAE negative ELBO
\end{minipage}
\hfill
\begin{minipage}{0.48\linewidth}
\centering
\includegraphics[width=\linewidth]{figures/ddpm_noise_prediction_loss.png}
\small (b) DDPM noise-prediction MSE
\end{minipage}
\caption{Training curves for the two models. The absolute loss values are not directly comparable because the objectives are different.}
\end{figure}

\end{document}
